\begin{textsample}{NEG Dim 4 – 2011\_c6 – Score -2.00 – 2011\_c6/enade\_2011\_eng\_g\_5\_6b7bd07019.txt}  \label{ex:f4_neg_058}
CAPÍTULO 3 ANÁLISE TÉCNICA DA PROVA Este capítulo tem por objetivo apresentar o desempenho dos estudantes concluintes de Engenharia - Grupo V no ENADE/2011.
Para isso, foram calculadas as estatísticas básicas da prova em seu todo, bem como as estatísticas dos componentes relacionadas à Formação Geral, ao de Conhecimento Específico da Área e das questões discursivas isoladamente.
Nas tabelas, são apresentados o tamanho da população inscrita e de presentes, e as seguintes estatísticas das notas8 : média do desempenho na prova, erro padrão da média, desvio padrão, nota mínima, mediana e nota máxima.
As estatísticas apresentadas neste capítulo contemplam o total de estudantes concluintes da área de Engenharia - Grupo V em 2011 do Brasil e, separadamente, por Grande Região.
Foram calculadas tendo -se em vista as seguintes agregações : ( a ) as Grandes Regiões e o país como um todo ; ( b ) a Categoria Administrativa ; e ( c ) a Organização Acadêmica.
Em relação aos gráficos de distribuição de notas, o intervalo considerado foi de 10 unidades, aberto à esquerda e fechado à direita, com exceção do primeiro intervalo, [ 0 ; 10 ], fechado em ambos os extremos.
Para os gráficos de distribuição das notas das questões discursivas, foram consideradas mais duas categorias : questão em branco e nota zero. 3.1 ESTATÍSTICAS BÁSICAS DA PROVA 3.1.1 Estatísticas Básicas Gerais A Tabela 3.1 apresenta as estatísticas básicas da prova por grande Região.
A população total de inscritos foi de 992.
Destes, 920 estiveram presentes, sendo 7,3 \% o índice de não comparecimento.
A Região de maior abstenção foi a Sudeste ( 8,0 \% ) e a de menor abstenção foi a Norte ( 5,0 \% ).
A região Centro-Oeste não apresentou inscritos. 8 Essas estatísticas e outras estão definidas no Capítulo 1.
A média das notas da prova como um todo ( nas seções seguintes serão analisados os componentes de Formação Geral e de Conhecimento Específico ) foi 52,1, sendo que os alunos da região Norte obtiveram a média mais baixa ( 46,5 ) e os da região Sudeste obtiveram a média mais alta ( 53,4 ).
As demais médias foram : 47,3 na região Nordeste e 51,1 na região Sul.
O desvio padrão para o Brasil como um todo foi 13,9, sendo o maior desvio padrão encontrado na região Sul ( 14,6 ) e o menor na região Nordeste ( 10,9 ), indicando uma menor dispersão das notas desta última região.
Não existem alunos inscritos desta Área na região Coentro-Oeste e consequentemente nenhuma estatística sobre a prova.
A região que obteve a maior nota máxima foi a Sudeste ( 88,5 ), ao passo que a região que atingiu a menor nota máxima foi a Nordeste ( 76,5 ).
A mediana do Brasil como um todo foi 52,6, sendo a maior mediana obtida na região Sudeste ( 54,2 ) e a menor obtida na Norte ( 44,7 ).
A nota mínima zero ocorreu apenas na região Sudeste.
As demais foram : 25,0 na região Norte, 24,4 na região Nordeste e 15,1 na região Sul.
O comportamento das notas dos estudantes de todo o Brasil pode ser observado no Gráfico 3.1 que apresenta um histograma com a distribuição das mesmas.
Essa é uma distribuição unimodal com moda no intervalo ( 50 ; 60 ].
Apesar do coeficiente de assimetria da distribuição das notas ser negativo ( -0,16 ) este é pequeno e podemos considerar que a distribuição é aproximadamente simétrica.
As distribuições por Grande Região também apresentam tanto assimetria negativa quanto positiva : as regiões Norte e Nordeste, com coeficientes de assimetria positivo ( 0,76 e 0,18, respectivamente ) com concentração à esquerda e cauda um pouco mais longa à direita ; e as regiões Sudeste e Sul, com coeficiente de assimetria negativo ( -0,30 e -0,09, respectivamente ) apresenta o comportamento oposto.
Os Gráficos 3.2, 3.3 e 3.4 apresentam informações referentes à média da nota final dos Participantes, desagregando os resultados de acordo com, respectivamente, as Grandes Regiões do país, a Categoria Administrativa e a Organização Acadêmica.
Os gráficos apresentam o valor da média das notas como uma barra e os extremos do intervalo de confiança de 95 \% como linhas verticais unidas por uma linha horizontal na forma da letra H maiúscula.
Considerando -se o gráfico de notas segundo Grande Região ( Gráfico 3.2 ), observa -se que existe diferença estatisticamente significativa ao nível de 95 \% entre a maior média, obtida na região Sudeste ( 53,4 ) e a menor, obtida na região Norte ( 46,5 ).
Levando -se em conta os estudantes por Categorias Administrativas ( Gráfico 3.3 ), observa -se que existe diferença estatisticamente significativa entre as médias das notas das IES Públicas ( 54,8 ) e Privadas ( 46,1 ).
Pode ser observado ainda que a média dos alunos de IES Públicas é maior que a dos alunos de IES Privadas.
A diferença entre IES Públicas e Privadas é maior do que as diferenças entre regiões.
Tendo como foco o Gráfico 3.4, que apresenta as notas médias das provas segundo Organização Acadêmica, constata -se que não existe diferença estatisticamente significativa ao nível de 95 \% nas médias das notas dos estudantes provenientes de Centros Universitários e Faculdades, mas na comparação deste com a Universidade passa a existir.
A maior média foi obtida pelos estudantes de Universidades ( 53,0 ), e a menor, pelos de Faculdades ( 45,1 ). 3.1.2 Estatísticas Básicas no Componente de Formação Geral A Tabela 3.2 apresenta as estatísticas básicas em relação ao componente da prova que avalia a Formação Geral dos estudantes concluintes.
Os alunos de todo Brasil obtiveram desempenho médio de 59,0.
Quanto à variabilidade, o desvio padrão das notas dos estudantes do Brasil como um todo foi 15,8.
A maior média foi obtida na região Sudeste ( 60,1 ), e a menor, na região Norte ( 54,5 ).
As demais médias foram : 57,0 nas regiões Nordeste e Sul.
Já o maior desvio padrão foi obtido na região Norte ( 18,1 ) e o menor na região Nordeste ( 14,0 ).
Os demais desvios padrões foram : 15,8 na região Sudeste e 16,0 na região Sul.
A maior nota no Componente de Formação Geral da prova do ENADE/2011 foi obtida por pelo menos um aluno da região Sul ( 97,0 ) enquanto que a menor nota máxima foi obtida na região Norte ( 82,5 ).
Nas outras regiões as notas máximas foram : 88,5 na região Nordeste e 92,5 na região Sudeste.
A mediana do Brasil como um todo foi 61,5, sendo a menor mediana encontrada na região Norte ( 54,0 ) e a maior encontrada na região Sudeste ( 62,0 ).
A nota mínima nesta parte foi zero nas regiões Sudeste e Sul, 15,0 na região Norte e 29,0 na região Nordeste.
O Gráfico 3.5 propicia a avaliação do desempenho dos estudantes no Componente de Formação Geral a partir do histograma da distribuição das notas correspondentes.
A distribuição é unimodal, com moda em ( 60 ; 70 ].
Nota -se, ainda, que no gráfico 3.5 as notas apresentam uma maior dispersão do que no Gráfico 3.1 ( distribuição das notas da prova ), confirmado pela comparação dos desvios padrões : 13,9 para a nota da prova como um todo e 15,8 para o Componente de Formação Geral.
Para o Componente de Formação Geral, o coeficiente de assimetria da distribuição das notas dos estudantes, como na prova como um todo, também é negativo ( -0,79 ).
A distribuição apresenta concentração à direita e cauda maior à esquerda.
Em todas as Grandes Regiões os histogramas também possuem assimetria negativa ( entre – 0,19 e – 0,92 ).
Nos Gráficos 3.6, 3.7 e 3.8 são apresentadas as informações referentes ao desempenho dos Concluintes no Componente de Formação Geral, em diferentes agregações : Grande Região do país, Categoria Administrativa e Organização Acadêmica.
Observa -se pelo Gráfico 3.6 que não existe diferença estatisticamente significativa entre as médias das notas no Componente de Formação Geral, segundo Grande Região do país.
Vemos que o intervalo de confiança mais largo é o da região Norte ; já o intervalo mais estreito é observado na região Sudeste.
Este fato está relacionado, também, com o tamanho da população envolvida, menor na região Norte do que na Sudeste.
No Gráfico 3.7, que representa as notas médias no Componente de Formação Geral segundo Categoria Administrativa do país, observa -se que existe diferença estatisticamente significativa entre as médias.
Semelhantemente ao que ocorreu para as notas da prova como um todo, para o Componente de Formação Geral, os concluintes de Engenharia - Grupo V, das IES Públicas ( 59,9 ) obtiveram uma média maior do que os das IES Privadas ( 56,9 ).
Considerando -se o tipo de Organização Acadêmica, nota -se, no Gráfico 3.8, não existe diferença estatisticamente significativa entre a maior e a menor média.
Nas Universidades ( 59,4 ), a média é mais elevada do que em Centros Universitários ( 56,7 ) e Faculdades ( 56,3 ). 3.1.3 Estatísticas Básicas do Componente de Conhecimento Específico A Tabela 3.3 apresenta as estatísticas básicas referentes ao Componente de Conhecimento Específico da área de Engenharia - Grupo V.
A média do desempenho dos alunos do Brasil como um todo foi 49,8.
A maior média foi obtida na região Sudeste ( 51,1 ), e a menor, na região Norte ( 43,8 ).
As demais médias foram : 44,0 na região Nordeste e 49,1 na região Sul.
Quanto à variabilidade das notas, o desvio padrão do Brasil como um todo foi 16,5, sendo o maior desvio padrão observado na região Sul ( 17,5 ) e o menor na região Nordeste ( 12,5 ).
Os demais desvios foram : 12,9 da região Norte e 16,6 da região Sudeste.
A mediana das notas dos estudantes de todo o Brasil foi 49,3.
A maior mediana ocorreu na região Sudeste ( 51,3 ) e a menor na região Norte ( 42,6 ).
As demais medianas foram 44,4 na região Nordeste e 49,1 na região Sul.
A nota máxima do Brasil como um todo foi 90,9, sendo obtida por pelo menos um aluno da região Sul.
As demais notas máximas foram : 80,3 na região Norte, 78,8 na região Nordeste e 90,7 na região Sudeste.
A nota mínima foi zero apenas na região Sudeste.
As demais notas foram : 22,2 na região Norte, 19,9 na região Nordeste e 12,2 na região Sul.
Assim como os Gráficos 3.1 e 3.5, o Gráfico 3.9, apresentado a seguir, proporciona uma avaliação do desempenho de concluintes em relação ao Componente de Conhecimento Específico com um histograma da distribuição das notas correspondentes.
Dentre as três distribuições apresentadas, esta é a mais concentrada nas notas baixas.
Esta também é uma distribuição unimodal, e o grupo modal é o ( 40 ; 50 ].
O coeficiente de assimetria da distribuição das notas do Componente de Conhecimento Específico é negativo e bem próximo de zero ( – 0,01 ).
Nota -se pelo histograma ( Gráfico 3.9 ) que esta é uma distribuição aproximadamente simétrica.
Já na região Norte, o coeficiente de assimetria é positivo, igual a 0,75, evidenciando que a distribuição das notas dos estudantes dessa região, no Componente de Conhecimento Específico, tem cauda mais pesada à direita.
Os Gráficos 3.10, 3.11 e 3.12 apresentam uma comparação dos resultados em relação à Grande Região do país, à Categoria Administrativa e à Organização Acadêmica, agora levando em conta o desempenho de alunos no Componente de Conhecimento Específico da prova.
Pelo Gráfico 3.10, observa -se que não existe diferença estatisticamente significativa entre as médias das notas no componente de Conhecimento Específico, das regiões Norte e Nordeste em relação à região Sul.
As diferenças são significativas com respeito à região Sudeste.
A região Sudeste ( 51,1 ) foi a que obteve média mais elevada, enquanto a região Norte ( 43,8 ) foi a região que obteve média mais baixa.
Quanto à Categoria Administrativa ( gráfico 3.11 ), observa -se um comportamento semelhante àquele da parte de Formação Geral e à prova como um todo, ou seja, existe diferença estatisticamente significativa entre as médias das IES Públicas ( 53,1 ) e Privadas ( 42,5 ), sendo que a maior média foi obtida por alunos de IES Públicas de ensino.
Quanto ao Gráfico 3.12, observa -se, mais uma vez, que existe diferença estatisticamente significativa ao nível de 95 \% entre as notas no Componente de Conhecimento Específico das Universidades em relação aos Centros Universitários e Faculdades.
A média dos concluintes das Universidades ( 50,8 ) foi maior do que de Centros Universitários ( 45,3 ) e de Faculdades ( 41,4 ). 3.2 ANÁLISE DAS QUESTÕES OBJETIVAS 3.2.1 Componente de Formação Geral A Tabela 3.4 apresenta as estatísticas básicas relativas às oito questões objetivas do componente da prova que abrange a Formação Geral dos estudantes.
A média do Brasil foi 55,6.
A menor média foi encontrada na região Nordeste ( 52,5 ) e a maior na região Sudeste ( 56,7 ).
As demais médias foram 52,6 na região Norte e 56,7 na região Sudeste.
O desvio padrão do Brasil foi 17,5, sendo o maior desvio padrão encontrado na região Nordeste ( 18,4 ) e o menor na região Sul ( 16,8 ).
Os demais desvios foram : 18,2 na região Norte e 17,5 na região Sudeste.
A mediana do Brasil é de 62,5 que se repete na região Sudeste, as demais regiões apresentam mediana 50,0.
As notas máximas de todas as regiões foi 100,0, com exceção a região Norte com 87,5.
As notas mínimas das regiões Sudeste e Sul foi 0,0, a da região Norte foi 25,0 e da região Nordeste foi 12,5.
Considerando -se a distribuição das notas obtidas nas questões objetivas do Componente de Formação Geral, o coeficiente de assimetria é sempre negativo, mesmo separando -se os estudantes por Grande Região, variando entre – 0,30 e – 0,40.
Isto significa que há uma cauda um pouco maior do lado esquerdo da distribuição.
A Tabela 3.5 apresenta os índices de facilidade e discriminação ( ponto bisserial ) das questões objetivas do Componente de Formação Geral.
Quanto ao índice de facilidade, foram usadas as seguintes cores para diferenciar o nível de dificuldade da questão : • Azul para as questões classificadas com índice muito fácil ( índice \&gt ; =0,86 ), verde para as questões classificadas com índice fácil ( 0,61 a 0,85 ), amarelo para as questões classificadas com médio ( 0,41 a 0,60 ), vermelho para as questões classificadas com difícil ( 0,16 a 0,40 ) e roxo para as questões classificadas com muito difícil ( \&lt ; =0,15 ).
Já quanto ao índice de discriminação, foram usadas as seguintes cores para qualificar a questão : • As questões classificadas com índice fraco receberam a cor vermelho ( índice \&lt ; =0,19 ), as classificadas com médio receberam a cor amarelo ( 0,20 a 0,29 ), as classificadas com bom receberam a cor verde ( 0,30 a 0,39 ) e as classificadas com muito bom ( \&gt ; =0,40 ) receberam a cor azul.
As questões objetivas do Componente de Formação Geral, segundo o índice de facilidade, foram assim avaliadas : das oito questões, nenhuma teve o índice de facilidade classificado como muito fácil.
Quatro questões foram tidas como fáceis, por terem índice de acertos situado na faixa entre 0,61 e 0,85 ( de 61,0 \% a 85,0 \% de acertos ).
Uma questão foi considerada de dificuldade média, situando -se no intervalo entre 0,41 e 0,60 do índice de facilidade, ou seja, houve entre 41,0 \% e 60,0 \% de acertos.
Duas questões foram classificadas como difícil, situando -se no intervalo entre 0,16 e 0,40.
Por fim, uma questão apresentou menos de 15 \% de acertos, razão pela qual foi classificada como muito difícil.
Como já comentado, para análise das questões objetivas relativas à Formação Geral segundo o poder de discriminação, utilizou -se, o índice de discriminação ponto bisserial.
Nesta análise as questões foram assim avaliadas : quatro das oito questões apresentaram índices acima de 0,40 e, assim, foram classificadas com índice muito bom para esse grupo de alunos ; três questões tiveram bom índice de discriminação, entre 0,30 e 0,39.
Uma questão foi classificada com índice médio.
O índice de facilidade variou de 0,15 a 0,82, e o de discriminação, de 0,26 a 0,46.
As questões com índices de discriminação muito bom, de números 1, 2, 3, e 5, figuraram entre as mais fáceis desse conjunto : três classificadas na categoria fácil ( questões 1, 3 e 5 ) do índice de facilidade e uma na categoria médio ( questão 2 ).
Em particular, as questões 2 e 3 foram as que apresentaram maior poder discriminatório, com índice 0,46, sendo a questão 3 a mais fácil com uma proporção de 0,82 de acertos.
A questão de número 8 apresentou índice de facilidade 0,15, ou seja, um quantitativo pequeno de estudantes conseguiu resolvê -la, dentro do universo de participantes.
Mas, seu índice de discriminação foi médio, o que não a eliminou do cômputo da nota final.
Quatro questões apresentaram índice de facilidade fácil ( 1, 3, 5 e 6 ) variando entre 0,73 a 0,82 e três delas ( 1, 3 e 5 ) com um índice de discriminação considerado muito bom, variando de 0,44 a 0,46 e uma ( questão 6 ) com índice de discriminação considerado bom ( 0,36 ).
O Gráfico 3.13, para exemplificar, analisa o comportamento da questão de número 3 de Formação Geral.
Trata -se da questão mais fácil e uma das que obteve o maior índice de discriminação dessa parte da prova.
Neste gráfico, cada uma das cinco curvas representa o percentual de respostas em determinada alternativa da questão, em função da nota dos estudantes nesta parte da prova ( Formação Geral/Múltipla Escolha ), antes de possíveis eliminações pelo critério do ponto bisserial.
A curva em vermelho corresponde à alternativa E, a correta para esta questão.
Assim, observa -se que entre os estudantes com notas mais baixas, nessa parte do exame, a situação mais frequente foi a escolha de uma das alternativas incorretas : a alternativa C ( em preto ) ou D ( em roxo ).
Na medida em que a nota aumenta, indicando desempenho melhor nesta parte da prova, aumenta concomitantemente a proporção de estudantes que selecionaram a alternativa correta B, atingindo 100 \% para sete acertos.
Essa análise permite verificar como a questão discriminou os grupos de desempenho, justificando o alto índice obtido na questão.
Os gráficos relativos às demais questões de Formação Geral constam do Anexo I. 3.2.2 Componente de Conhecimento Específico A Tabela 3.6 apresenta as estatísticas básicas em relação às questões objetivas do Componente de Conhecimento Específico da prova, por Grande Região.
A média do Brasil deste componente foi de 54,3.
A menor média foi observada na região Norte ( 47,8 ) e a maior na região Sudeste ( 55,7 ).
O desvio padrão de todo o Brasil foi 17,6, sendo o menor desvio padrão encontrado nas regiões Norte e Nordeste ( 13,4 ) e o maior na região Sul ( 18,5 ).
A mediana de todo o Brasil foi 54,5.
A menor mediana foi observada na região Norte ( 46,6 ) e a maior na região Sudeste ( 55,0 ).
A nota máxima da prova ( 95,7 ) foi obtida no componente de Conhecimento específico por pelo menos um aluno das regiões Sudeste e Sul.
As regiões Norte e Nordeste apresentaram as menores notas máximas, ambas com 85,7.
Somente a região Sudeste apresentou nota mínima zero.
A Tabela 3.7 apresenta os índices de facilidade e discriminação ( ponto bisserial ) das questões objetivas do Componente de Conhecimento Específico.
Para facilitar a diferenciação das questões usou -se as mesmas cores da Tabela 3.5 para as diferentes classificações dos índices de facilidade e de discriminação.
Dentre as questões objetivas da parte da prova relativa ao Componente de Conhecimento Específico, as de número 29 e 34 foram anuladas pela Comissão.
Desse modo, a classificação quanto ao índice de facilidade foi estabelecida com base em 36 das 38 questões.
A partir dos índices obtidos, pode -se concluir que a classificação modal das questões objetivas da prova com respeito à dificuldade foi a média : das 36 questões, dezesseis foram classificadas como média.
Não houve questão classificada como muito fácil.
Treze questões foram tidas como fáceis, na faixa de 0,62 a 0,78 do índice de facilidade, seis como difícil e uma como muito difícil.
Já quanto aos índices de discriminação das questões objetivas do Componente de Conhecimento Específico da prova, tem -se como resultado a seguinte classificação : treze das 36 questões válidas foram consideradas como boas, enquanto dez delas tiveram índice de discriminação muito bom.
Assim, para mais da metade das questões – 23 em 36 – os índices de discriminação foram bons ou muito bons.
Dentre as demais, sete delas foram classificadas como médias e outras seis como fracas, sendo treze, por conseguinte, a quantidade de questões nos dois patamares mais baixos de discriminação.
Constata -se, assim, que a prova – no que se refere ao Componente de Conhecimento Específico – possuía capacidade média de discriminar entre aqueles que dominam ou não o conteúdo.
Dentre as questões que alcançaram os maiores índices de discriminação, as de números 15, 22, 24, 25, 27, 33, 39, 42, 44 e 46 classificadas com índice muito bom, situando -se no intervalo de 0,40 a 0,55 do índice, cinco delas ( questões 24, 27, 39, 44 e 46 ) foram classificadas na categoria fácil, quatro ( 22, 25, 33 e 42 ) como média e uma ( 15 ) como difícil, quanto ao índice de facilidade.
A questão de número 10 foi a mais difícil dentre as 36 questões específicas válidas, com baixo índice de facilidade, apenas 12,0 \% de acertos.
Essa questão apresentou poder discriminatório igualmente baixo, 0,16, o que comprova ter sido esta a mais difícil para os estudantes.
Destaca -se, também, a questão 11, com índice de facilidade 0,25, o que, em termos percentuais, corresponde a 25,0 \% de estudantes que responderam acertadamente, obtendo, ainda, 0,09 de índice de discriminação.
Tais questões foram, portanto, duas das mais difíceis da prova.
Por isso, as questões 10 e 11 foram eliminadas do cômputo da nota final.
Além destas duas, as demais questões com índice fraco de discriminação, questões 14, 31, 43 e 45 também não foram computadas.
A título de exemplo das análises do comportamento das questões objetivas, o Gráfico 3.14 analisa a questão 44 do Componente de Conhecimento Específico.
Esta foi uma das questões mais fáceis da prova, apresentando índice de facilidade 0,63, ou seja, 63,0 \% dos estudantes assinalaram acertadamente a opção C, correspondente ao gabarito.
Seu índice de discriminação foi igual a 0,55, classificado como muito bom.
Neste gráfico, cada uma das cinco curvas representa o percentual de respostas em determinada alternativa da questão 44, em função da nota dos estudantes nesta parte da prova, antes de possíveis eliminações de questões pelo critério do ponto bisserial.
A alternativa correta C, representada no gráfico pela curva em preto, foi escolhida em maiores proporções pelos alunos com desempenho melhor nesta parte da prova.
Já as alternativas incorretas, também denominadas distratores, foram selecionadas principalmente por aqueles com notas mais baixas.
Aqueles com nota zero, na sua quase totalidade deixaram esta questão em branco ou marcaram mais de uma alternativa, comportamento considerado inválido.
A proporção de alunos que selecionou a resposta correta C aumenta gradativamente, chegando a atingir 100 \% para os estudantes com 24 acertos, enquanto a proporção dos que escolheram alternativas incorretas decai, a partir de quatro acertos, como função da nota nesta parte da prova.
Os gráficos relativos às demais questões do Componente de Conhecimento Específico constam do Anexo I. 3.3 ANÁLISE DAS QUESTÕES DISCURSIVAS 3.3.1 Componente de Formação Geral As análises dos resultados de desempenho dos estudantes de Engenharia - Grupo V nas duas questões discursivas relativas à Formação Geral encontram -se na Tabela 3.8 e no Gráfico 3.15.
Na tabela 3.8 observa -se que as notas médias foram mais baixas nesse conjunto de questões do que no das objetivas.
Os estudantes de todo o Brasil obtiveram, em Formação Geral, média 55,6 nas questões objetivas e 64,0 nas questões discursivas.
A mediana confirma o melhor desempenho dos alunos de todo o Brasil nas questões discursivas do componente de Formação Geral.
Enquanto essa estatística foi de 70,0 para questões discursivas, para as questões objetivas essa estatística foi de 62,5.
Pode -se, também, notar um aumento do desvio padrão de 17,5, nas questões objetivas do Componente de Formação Geral dos alunos de todo o Brasil, para 25,4 nas questões discursivas do mesmo componente.
Como já comentado, a mediana de todo o Brasil, neste componente, foi 70,0, sendo a maior mediana encontrada na região Sudeste ( 72,5 ) e a menor na região Norte ( 62,5 ).
A nota máxima ( 100,0 ) e a mínima ( 0,0 ) foram obtidas por alunos de todas as regiões do Brasil.
O Gráfico 3.15 representa a distribuição das notas nas questões discursivas no Componente de Formação Geral.
A moda desta distribuição ocorre no intervalo ( 70 ; 80 ], devido à grande quantidade de notas elevadas.
Nota -se um máximo local no primeiro intervalo [ 0 ; 10 ].
Na sequência, os resultados verificados para cada uma das questões discursivas de Formação Geral serão apresentados, estabelecendo -se relações com os conteúdos abordados em cada uma delas.
Os comentários da Banca de docentes corretores a respeito do observado na correção das respostas dos estudantes, suas impressões e conclusões serão apresentados junto à análise de cada questão.
Cumpre esclarecer que, tendo em vista que as questões discursivas de Formação Geral são padronizadas, ou seja, constam de todas as provas, os comentários da Banca são os mesmos para todas as carreiras acadêmicas, sendo direcionados a todos os estudantes que participaram do ENADE/2011.
A seguir, serão analisados os desempenhos da Área de Engenharia - Grupo V nas duas questões discursivas de Formação Geral do ENADE/2011, comparando os resultados obtidos com comentários para cada questão. 3.3.1.1 Análise da Questão Discursiva 1 do Componente de Formação Geral Os dados de Engenharia - Grupo V, obtidos a partir das respostas à questão 1, encontram -se na Tabela 3.9 e no Gráfico 3.16.
Nessa questão – de melhor desempenho dentre as duas de Formação Geral – os alunos de todo Brasil tiveram média, 71,4.
A maior média para a questão 1 foi obtida na região Sudeste ( 73,1 ), e a menor, na região Norte ( 59,6 ).
Quanto à variabilidade das notas, o desvio padrão de todo o Brasil foi 31,3.
O menor desvio padrão foi obtido na região Nordeste ( 28,9 ) e o maior desvio padrão foi obtido na região Norte ( 37,1 ).
As notas máximas e as mínimas da questão discursiva 1 foram as mesmas para todas as regiões do Brasil, respectivamente, 0,0 e 100,0.
A mediana foi menor nas regiões Norte e Nordeste ( 70,0 ) e maior nas regiões Sudeste e Sul ( 90,0 ).
O Gráfico 3.16 mostra a distribuição das notas na questão discursiva 1 do Componente de Formação Geral.
Observa -se que a maior frequência corresponde aos alunos que tiraram notas no último intervalo, sendo este o intervalo modal.
Observa -se também uma grande concentração de notas no intervalo ( 80 ; 90 ].
Um pouco mais de 5 \% dos participantes deixaram esta questão em branco.
O índice de assimetria foi -0,89. 3.3.1.2 Comentários sobre a correção das respostas à Questão Discursiva 1 De maneira geral, a aplicação da chave de correção da questão 1 de Formação Geral não apresentou qualquer dificuldade digna de menção.
Para isso contribuíram, indubitavelmente, os ajustes feitos após a correção da amostra e a reunião entre todos os membros da Banca de docentes corretores.
As poucas dúvidas, todas pontuais, apresentadas pelos corretores, foram acompanhadas e respondidas pela coordenação e subcoordenação da correção das questões de Formação Geral, por meio da ferramenta de Gerenciamento de Dúvidas do Sistema de Correção On-line.
Não houve registro de qualquer ocorrência que pusesse em xeque o padrão de resposta ou a efetividade e a adequação da chave de correção.
Explica -se : trata -se de questão com comando claro, direto e objetivo ( solicitava -se, basicamente, três vantagens justificadas de cursos a distância ), cujas respostas foram corrigidas por meio da aplicação de um chave de correção testada e aprovada previamente.
Havia absoluta clareza quanto aos critérios de avaliação da correspondência entre as respostas dos estudantes e as possibilidades de vantagens de cursos a distância admitidas como corretas no padrão de resposta oficial, além de gradações explícitas ( e fáceis de aplicar ) dos diferentes níveis de pontuação previstos.
Felizmente, portanto, não há \textbf{reparo} a registrar em relação à facilidade de aplicação do padrão de resposta e da chave de correção, e nem em relação à atribuição dos diferentes níveis de pontuação previstos.
Todas as dificuldades que poderiam ter obstado a correta aplicação do padrão de resposta oficial e da respectiva chave de correção foram evitadas por meio dos ajustes feitos após a correção da amostra e fartamente debatidos com toda a Banca.
Digno de nota é que quantidade tão significativa de profissionais envolvidos na mesma tarefa – tanto para a questão 1 quanto para a 2 – tenha apresentado tão poucas dificuldades na execução da correção, em termos proporcionais.
Em suma, a correção da questão 1 da prova de Formação Geral do ENADE 2011 foi exemplarmente bem planejada, servindo -se de padrão de resposta muito bem adequado à questão proposta.
Quanto ao tema desta questão, em particular, a Banca verificou que uma parcela significativa de estudantes evocou experiências bastante concretas e próximas de sua realidade.
Houve várias respostas que indicavam uma vivência pessoal de ensino superior na modalidade Educação a Distância ( EaD ), evidenciando o tom de depoentes nos textos apresentados.
Os estudantes, em sua grande maioria, utilizaram parte considerável do espaço de 15 linhas disponíveis para a resposta – e outra parcela menos significativa dos que não o fizeram demonstraram notável capacidade de atender ao comando da questão de maneira objetiva, curta e, via de regra, correta.
Registre -se, ainda, que foram relativamente poucos os casos de respostas que tenham passado ao largo do tema em pauta na questão.
As capacidades de leitura, de compreensão do comando proposto e de expressão escrita que os estudantes avaliados na edição 2011 do ENADE foram satisfatórias.
Não obstante, seria leviano perder de vista que a qualidade dos textos redigidos em resposta às questões discursivas do Exame ainda está muito aquém do que se espera de concluintes de cursos de ensino superior de todas as regiões do país.
Quanto ao conteúdo das respostas, a Banca constatou boa capacidade, por parte da maioria dos estudantes, de compreensão do tema e do comando da questão.
Foram relativamente poucos os casos de respostas que deixaram de enumerar vantagens da modalidade EaD, e proporcionalmente escassos os estudantes que citaram vantagens não previstas no padrão de resposta.
Foi frequente, a tentativa direta de atender ao comando da questão.
Os erros mais comuns, em relação ao padrão de respostas e à grade de correção, foram fruto do desdobramento em vários “ itens ” daquilo que, de acordo com o padrão de respostas oficial, representava uma única vantagem.
Destacaram -se, neste caso, as respostas que apontavam a flexibilidade de horário e/ou local como duas vantagens distintas daquela modalidade de ensino.
Quanto aos diferentes níveis de pontuação previstos, a maior causa de baixas pontuações foi a ausência de justificativas, e mesmo de argumentação, para uma ou mais das vantagens enumeradas.
Isso demonstra não apenas a objetividade das respostas, coerente com a objetividade do comando da questão ( “ enumere três vantagens de um curso a distância ” ), mas também certa dificuldade de formulação plena de um texto, ou ao menos de parágrafos, em formato dissertativo – mesmo diante de uma média de 5 linhas disponíveis para cada vantagem a enumerar.
Os acertos mais comuns, ou seja, os “ itens ” do padrão de respostas mais frequentemente mencionados foram : ( 1 ) a flexibilidade de horário e/ou local ; ( 2 ) a capilaridade do ensino a distância ; ( 3 ) a democratização do acesso à educação de qualidade ; e ( 4 ) os custos menores que os de cursos presenciais.
Dentre as vantagens previstas no padrão de respostas que foram menos citadas, destacam -se a inclusão de pessoas com comprometimento motor, a qualificação de professores e a troca de experiências entre os participantes.
Foram muito comuns, no entanto, as menções à supostamente maior facilidade de acesso a professores ou/e tutores em cursos superiores a distância. 3.3.1.3 Análise da Questão Discursiva 2 do Componente de Formação Geral A Tabela 3.10 mostra que o desempenho dos estudantes na questão 2 ( média 56,6 ) foi inferior ao obtido na questão de número 1 ( média 71,4 ).
A região Nordeste foi aquela onde a média, nessa questão, foi maior ( 58,5 ), e a de menor média foi a região Sul ( 53,2 ).
Quanto à variabilidade das notas, o desvio padrão de todo o Brasil foi 29,8, inferior ao obtido na questão de número 1 ( 31,3 ).
O maior desvio nessa questão foi obtido na região Norte ( 33,0 ), enquanto o menor foi obtido na região Sul ( 29,3 ).
A mediana comprova a diferença de desempenho entre a região Nordeste e as demais regiões ( 60,0 ), apresentando uma mediana de valor superior ( 70,0 ).
A nota máxima ( 100,0 ) e a mínima ( 0,0 ) foram obtidas em todas as regiões do Brasil.
O Gráfico 3.17 mostra a distribuição das notas na questão discursiva 2 do Componente de Formação Geral.
Este gráfico apresenta uma distribuição bimodal.
Sendo a moda principal na união dos intervalos ( 50 ; 60 ] e ( 60 ; 70 ].
A moda secundária, intervalo com a terceira maior frequência, é aquele que representa as questões deixadas em branco ( um pouco mais de 13 \% ).
Nota -se, ainda, que as notas ficaram distribuídas ao longo dos intervalos, a partir do ( 10 ; 20 ].
Diferente da questão discursiva de número 1 que concentrou as notas, principalmente, nos dois últimos intervalos.
O coeficiente de assimetria foi -0,70. 3.3.1.4 Comentários sobre a correção das respostas à Questão Discursiva 2 Diferentemente da questão 1, a questão 2, cujo tema aborda políticas públicas para a erradicação do analfabetismo e para o aumento da empregabilidade, permitiu mais amplitude na elaboração das respostas dos estudantes, uma vez que solicita a apresentação de uma proposta.
À falta de objetividade técnica na elaboração da questão – sobretudo no seu comando – correspondeu uma miríade de possibilidades interpretativas.
Registre -se que não raras foram as respostas que se utilizaram de lugares comuns e exortações religiosas/humanitárias/cívicas, contudo, na grande maioria dos casos, essas exortações foram usadas como um \textbf{encerramento} do texto e não comprometeram a resposta.
Muito frequente também foi a confusão entre política pública e política partidária, bem como alguns poucos, confundiam programas educacionais com programas educativos nas TVs.
Assim, destacaram -se como propostas/programas mais recorrentes :  prosseguimento das políticas já existentes, mas com o aumento dos investimentos, normalmente sugerindo a ampliação da bolsa família, da Educação de Jovens e Adultos ( EJA ) e da EaD ;  valorização do magistério, construção de escolas e melhoria das atuais ;  parceria do governo com empresas para manutenção de salas de aula visando a alfabetização de seus funcionários, oferecimento de estágios e redução da carga horária em troca de isenção de impostos ;  parcerias com igrejas e ONGs para criação de espaços de alfabetização ;  escolas itinerantes e alfabetizadores em domicílio, principalmente para pessoas com dificuldade de locomoção, como os idosos e deficientes físicos ;  erradicação do trabalho infantil ;  vinculação da bolsa família não apenas à frequência, mas também e, principalmente, ao resultado obtido pelo aluno na escola ;  revisão das políticas atuais, sendo a mais frequente a extinção da aprovação automática ;  críticas consistentes em relação a modelos didáticos considerados inadequados e desestimulantes para a educação não só de adultos, mas de pessoas de todas as idades.
Destacam -se, ainda, outras sugestões apresentadas :  Creches nas escolas onde os pais estão sendo alfabetizados.  Diminuição de duas horas na jornada de trabalho em empresas para funcionários não alfabetizados, para que possam frequentar a escola.  Campanhas educativas vinculadas aos meios de comunicação.  Aumento do número de escolas noturnas.  Formação específica para professores alfabetizadores.
Quanto à relação entre o analfabetismo e a empregabilidade, deve -se sublinhar que nem todos os estudantes estabeleceram claramente o vínculo entre essas duas situações sociais.
Alguns falaram separadamente de uma e de outra.
Mas a maioria fez referência à necessidade de estudo para “ conseguir um bom emprego com um bom salário ”.
Alguns, em menor número, estabeleceram de forma bastante interessante a questão histórica para a situação do Nordeste ; e, também, a relação entre a pessoa analfabeta/com pouca instrução e a desempregada/empregada em serviços mais pesados e pior remunerados, que não teriam condições de educar seus filhos que, por sua vez, também não teriam melhores oportunidades no mercado de trabalho, identificando a formação de um círculo vicioso e a necessidade de sua interrupção.
Muitos afirmaram que o analfabetismo não é o único responsável pelo desemprego e sim a má distribuição de renda.
Grande também foi o número de alunos que criticaram o resultado da pesquisa, afirmando que a mesma é enganosa, uma vez que considera alfabetizada a pessoa que “ desenha seu nome ”, sendo comum associarem a atual política de alfabetização com ganhos eleitorais.
Foi comum a resposta incompleta, em que o estudante fez apenas a análise das desigualdades/crítica do quadro apresentado, ou só apresentou proposta.
Alguns se limitaram a responsabilizar o governo referindo -se de forma bastante genérica à questão da “ educação ” e “ profissional ”.
Entre os equívocos que mais se repetiram, destacou -se a simples análise dos dados apresentados na tabela que consta do enunciado.
Alguns textos, inclusive, estavam corretos, sem que, no entanto, fosse respondida a pergunta.
Da mesma forma, quando os estudantes partiram para a segunda parte da questão, fizeram referência a vários programas já existentes ou simplesmente disseram que é muito importante que existam projetos para a educação.
Outro equívoco recorrente foi a análise da educação no Brasil como um todo.
Essa análise, apesar de correta, não se referia especificamente à questão do analfabetismo, o que, por conseguinte, levava a sugestões que não eram direcionadas à erradicação do analfabetismo, e sim à melhora da educação no Brasil.
Sendo assim, foram apresentadas propostas como o aumento do número de faculdades, o reforço de alunos do Ensino Médio, a criação de escolas técnicas, etc.
Por outro lado, houve um grande número de redações bem escritas e precisas, no que se refere ao que foi exigido pela questão.
Foram análises equilibradas e \textbf{sensatas}, correlacionando a problemática do analfabetismo com o desemprego e a apresentação de sugestões bastante consistentes. 3.3.2 Componente de Conhecimento Específico Na parte da prova relativa às questões discursivas no Componente de Conhecimento Específico ( Tabela 3.11 ), observa -se que a média foi mais baixa do que para as questões discursivas do Componente de Formação Geral.
Enquanto no Componente de Formação Geral a média para estudantes de Engenharia - Grupo V de todo o Brasil foi 64,0, na parte de Conhecimento Específico a média foi 24,1.
A maior média deste componente foi obtida pelos estudantes da região Sudeste ( 25,0 ), e a menor, pelos da região Norte ( 21,1 ).
Quanto à variabilidade das notas, o desvio padrão de todo o Brasil foi 19,2.
O maior desvio padrão foi encontrado na região Sul ( 20,1 ), e o menor, na região Nordeste ( 16,5 ).
A maior nota máxima foi obtida na região Sudeste ( 81,7 ) enquanto a menor nota máxima foi encontrada na região Nordeste ( 61,7 ).
Além disso, a nota mínima ( 0,0 ) foi obtida por alunos de todas as regiões do Brasil sem exceção.
A mediana do Brasil como um todo foi 21,7.
As medianas foram iguais entre si para as regiões Norte, Nordeste e Sul ( 20,0 ) e a região Sudeste apresentou uma mediana mais elevada ( 23,3 ).
O Gráfico 3.18 representa a distribuição das notas nas questões discursivas no Componente de Conhecimento Específico.
Essa distribuição é unimodal com moda no intervalo de ( 0 ; 10 ].
Nota -se uma tendência decrescente a partir do intervalo ( 10 ; 20 ] até o intervalo ( 80 ; 90 ].
Não foram registradas notas maiores que 90.
O coeficiente de assimetria foi 0,49.
A análise de cada uma destas questões será feita a seguir. 3.3.2.1 Análise da Questão Discursiva 3 do Componente de Conhecimento Específico Na questão 3, cujos resultados aferidos encontram -se descritos na Tabela 3.12, a média dos estudantes de todo o Brasil foi 19,6.
A menor média nessa questão foi obtida pelos alunos da região Nordeste ( 16,3 ), enquanto a maior média foi obtida na região Sudeste ( 20,5 ).
Quanto à variabilidade das notas, o desvio padrão de todo o Brasil foi 22,8.
O maior desvio padrão foi obtido na região Norte ( 24,2 ), enquanto o menor foi obtido na região Nordeste ( 21,0 ).
A nota máxima, 100,0 pontos, foi alcançada por pelo menos um aluno da região Sudeste.
As regiões onde ocorreram a menor nota máxima foram a Nordeste e a Sul ( 85,0 ).
A mediana do Brasil como um todo foi 10,0, a região Sudeste apresentou mediana superior ao Brasil ( 15,0 ) e as demais regiões abaixo : Norte e Sul ( 5,0 ) e Nordeste ( 7,5 ).
A nota mínima ( 0,0 ) foi obtida em todas as regiões do Brasil, sem exceção.
O Gráfico 3.19 mostra a distribuição das notas na questão discursiva 3, do Componente de Conhecimento Específico.
Essa distribuição é bimodal com moda nas questões em branco ( quase 35 \% ).
Analisando o gráfico sem considerar as questões em branco, a moda da distribuição está no intervalo ( 10 ; 20 ].
O índice de assimetria foi 1,02.
Em todas as regiões o índice de assimetria foi positivo e acima de 0,97 ( valor encontrado para a região Sudeste ), indicando distribuições semelhantes. 3.3.2.2 Comentários sobre as respostas à Questão Discursiva 3 A questão 3 abordou o tópico Processos de Soldagem.
Há vários processos de soldagem disponíveis, e cada um deles está associado a uma magnitude da Zona Afetada Pelo Calor ( ZAC ).
Um bom processo de soldagem é aquele que conduz a uma ZAC menor possível, já que a ZAC é uma região problemática do material, sujeita a defeitos, que enseja uma vulnerabilidade do mesmo.
A questão 3, no seu enunciado, mencionou quatro processos conhecidos de soldagem : Eletrodo revestido ; MIG ; MAG ; TIG.
A questão solicitou aos estudantes a redação de um texto, tendo como “ pano de fundo ” o Controle da ZAC, envolvendo 3 subitens : ( a ) ; ( b ) ; ( c ).
No subitem ( a ), foi solicitada a influência de diferentes processos de soldagem.
O padrão de resposta estabeleceu a abordagem de, no mínimo, 3 processos e seus reflexos na ZAC.
Cada processo apresentado pelos alunos correspondeu a 1,0 ponto para o total de 3,0 pontos deste subitem.
Alguns alunos, com muita propriedade, citaram outros processos importantes de soldagem não mencionados no enunciado, como por exemplo : Arco submerso ; Gás Oxi-acetileno ; Laser ; Feixe de elétrons ; etc.
A Banca valorizou tais respostas, de acordo com o padrão de resposta, a questão do aporte de energia característico de cada processo, bem como a Potência Específica por unidade de área da fonte de calor, afetam diretamente o tamanho da ZAC.
Quanto maior a potência específica por unidade de área, menor será a ZAC.
Logo a ZAC ( MIG/MAG/TIG ) será menor do que a ZAC do Eletrodo Revestido, que será menor do que a ZAC do eletrodo submerso, já que possuem potências específicas de 105, 104, e 103 W/m2, respectivamente.
A ZAC da soldagem a laser é bem pequena, menor do que MIG/MAG/TIG, e a soldagem a gás resulta em ZAC muito grande.
Apenas um reduzido número de alunos demonstrou conhecimento das ordens de grandeza de energia, dos diferentes processos de soldagem e as suas consequências sobre a ZAC.
No subitem ( b ), solicitaram -se os tipos de proteção da poça de fusão.
Sabe -se que o processo de soldagem exibe uma vulnerabilidade devido à exposição da poça de fusão a contaminações oriundas da atmosfera circunvizinha, como oxidação, introdução de elementos químicos indesejáveis deletérios, etc.
As proteções mais corriqueiras da poça de fusão são fornecidas usando -se atmosferas inertes ( processo MIG ), atmosferas de gases ativos ( processo MAG ), eletrodo revestido, cujo revestimento externo gera gases e uma camada de escória que protege a poça de fusão, arco submerso e seus fluxos de proteção.
Neste subitem, a Banca detectou um número maior de acertos em relação ao subitem anterior, por parte dos estudantes, já que os processos de soldagem apresentados no subitem ( a ) também foram analisados quanto a este aspecto de proteção.
Finalmente, no subitem ( c ), solicitou -se discutir aspectos relativos à microestrutura e macroestrutura, ou, alternativamente, os ensaios utilizados para analisá -los.
De acordo com o padrão de resposta, os efeitos na microestrutura a ZAC manifestam -se na alteração da estrutura dos grãos, nas fases presentes, na estrutura dendrítica de solidificação, e nos defeitos surgidos.
Quanto à macroestrutura, a junta soldada pode exibir trincas e poros, falta de preenchimento, entre outros defeitos.
Quanto aos ensaios, esperou -se que os estudantes mencionassem a realização de análises microestruturais por microscopia da região soldada e ensaios não destrutivos, como por exemplo, líquidos penetrantes e ultrassom, para fins de avaliação da qualidade metalúrgica da ZAC.
Neste subitem, tanto os aspectos micro e macroestruturais da junta soldada, quanto os ensaios nela realizados, foram considerados para fins de pontuação.
Fazendo uma análise resumida das respostas elaboradas pelos estudantes, foi possível perceber que o tema soldagem precisa ser mais abordado na formação dos alunos de Engenharia dos Materiais. 3.3.2.3 Análise da Questão Discursiva 4 do Componente de Conhecimento Específico A Tabela 3.13 contém as informações relativas à questão 4 do conjunto de questões do Componente de Conhecimento Específico.
O desempenho dos estudantes de todo o Brasil nesta questão foi superior ao desempenho na questão de número 3.
A média geral do Brasil foi 24,3, sendo a menor média registrada na região Norte ( 22,1 ) e a maior na região Sudeste ( 25,1 ).
A nota máxima ( 100,0 ) foi atingida na região Sudeste, as demais regiões foram : 80,0 na região Norte, 90,0 na região Nordeste e 85,0 na região Sul.
A mediana ( 10,0 ) foi a mesma em todas as regiões, a exceção da região Norte ( 5,0 ).
A nota mínima ( 0,0 ) foi obtida em todas as regiões do Brasil, sem exceção.
O Gráfico 3.20, representa a distribuição de notas na questão discursiva 4, no Componente de Conhecimento Específico.
Desconsiderando a categoria “ em branco ”, a distribuição é bimodal nos intervalos ( 0 ; 10 ] e ( 40 ; 50 ].
Observa -se que uma proporção de estudantes bem parecida, mas um pouco menor, deixou essa questão em branco ou recebeu nota zero, comparando -se à questão anterior.
Cerca de 18 \% dos alunos receberam nota no intervalo ( 40 ; 50 ] nesta questão. 3.3.2.4 Comentários sobre as respostas à Questão Discursiva 4 A questão foca na produção de grades de chumbo com uma produção diária muito alta.
Para isso, foram fornecidas duas tabelas.
A primeira exibia aspectos técnicos, com informações numéricas, e a segunda com aspectos econômicos qualitativos, referentes a 6 ( seis ) processos conhecidos de fundição.
A questão foi formulada por meio de um enunciado muito extenso, o que aparentemente confundiu os estudantes.
Revelou -se cansativa, o que se confirma pelo comentário crítico de uma estudante reproduzido fielmente a seguir. “ As características dos processos de fundição fornecidas estão confusas, não fornecendo informações suficientes para a definição do processo produtivo ”.
Esta questão foi dividida em subitens de respostas.
O subitem ( a ) referiu -se às características técnicas de cada processo mais adequadas à produção desejada ; o subitem ( b ) referiu -se às características econômicas de cada processo para o mesmo fim ; e o subitem ( c ) referiu -se à escolha do processo de fundição mais indicado para a produção desejada.
O padrão de resposta, estabelecido após a correção amostral, incluiu no processo decisório do melhor processo, à luz do aspecto “ custo x benefício ”, dois processos : a fundição em molde permanente e a fundição sobre pressão.
Isto fez com que a avaliação dos estudantes melhorasse significativamente, em relação à correção amostral.
Uma grande fração de alunos optou por escolher, como melhor processo de fundição para a produção das grades de chumbo, as moldagens em areia ou em gesso.
Tais participantes cometeram o equívoco de não atentar para a significativa quantidade de peças/dia a serem produzidas, e optaram pelo baixo custo de processos de produção lenta.
Tal abordagem exemplifica as respostas fracas ocorridas com maior frequência.
Como respostas medianas, alguns estudantes souberam selecionar os processos esperados de fundição, constantes do padrão de resposta, porém não apontaram a elevada produtividade desses processos, que era o fator determinante para esta escolha.
As respostas adequadas foram aquelas nas quais os participantes identificaram os processos de fundição sob pressão, e/ou em molde permanente, como aqueles mais indicados para a produção desejada, justificando -os acertadamente com base nas características técnicas e econômicas.
Comparando -se o resultado desta questão ( foco em fundição ), com a questão anterior ( foco em soldagem ), obteve -se uma média um pouco superior.
Quanto ao número de questões em branco, para esta questão também foi bastante alto, embora ligeiramente inferior do que o ocorrido na questão 3. 3.3.2.5 Análise da Questão Discursiva 5 do Componente de Conhecimento Específico A Tabela 3.14 contém as informações relativas à questão 5 do conjunto do Componente de Conhecimento Específico.
O desempenho dos estudantes nessa questão foi superior ao das questões 3 e 4.
A nota média dos estudantes de todo o Brasil foi 28,5.
A maior média foi registrada na região Sudeste ( 29,4 ), enquanto a menor média foi registrada na região Norte ( 23,4 ).
Quanto à variabilidade das notas, o desvio padrão dos alunos do A nota máxima ( 100,0 ) foi alcançada nas regiões Sudeste e Sul, enquanto a menor nota máxima foi registrada na região Norte ( 70,0 ).
A nota mínima foi zero para todas as regiões do Brasil, sem exceção.
A mediana do Brasil foi 20,0 e apenas a região Sudeste obteve valor acima deste ( 25,0 ), as demais regiões foram : 15,0 na região Norte, 20,0 na região Nordeste e 10,0 na região Sul.
O Gráfico 3.21 representa a distribuição das notas da questão discursiva 5 no Componente de Conhecimento Específico.
Destaca -se o grande número de estudantes que deixaram a questão 5 em branco ou obtiveram nota zero tendo respondido a questão.
No entanto, entre as três questões discursivas do Componente de Conhecimento Específico, esta foi a que apresentou a menor proporção de alunos que deixaram a questão em branco.
A frequencia da distribuição das notas é, grosso modo, decrescente com a nota.
O índice de assimetria para o Brasil foi 0,59.
Para as regiões, os índices foram da mesma ordem de grandeza, sugerindo distribuições com forma semelhante. 3.3.2.6 Comentários sobre as respostas à Questão Discursiva 5 A questão 5 explorou a propriedade mecânica de DUREZA, exibida pelas principais classes de materiais : metais ; cerâmicos e polímeros.
O padrão de resposta estabeleceu que a hierarquização pura e simples da magnitude da dureza : “ Dureza dos cerâmicos maior do que a dos metais e estes com dureza superior a dos polímeros ”, valeria com 40 \% da questão.
A partir daí, a justificativa para cada comportamento de dureza, de cada classe, ampliaria a pontuação.
De acordo com o padrão de resposta, todas as justificativas teriam que ter como ponto de partida o tipo de ligação atômica e molecular atuantes nos materiais.
Os estudantes que apontaram pelo menos este fato, mesmo errando na hierarquização e/ou justificativas, garantiram pelo menos 10 \% do valor total da questão.
Como respostas fracas, a Banca percebeu que alguns estudantes : • confundiram dureza com resistência mecânica ou mesmo tenacidade, e analisaram as outras propriedades mecânicas e não a dureza, escopo da questão ; • associaram o processo de fabricação dos citados materiais ( metais, cerâmicos e polímeros ) como sendo, erroneamente, responsáveis pela dureza dos mesmos. • conduziram a solução da questão, erroneamente, descrevendo técnicas de medição de dureza, evidenciando um flagrante erro de interpretação. • afirmaram que os polímeros exibiam menor dureza devido à presença das ligações covalentes, enquanto as ligações covalentes são ligações primárias fortes e não são elas que determinam a baixa dureza, e sim o elo mais fraco dos polímeros que são as ligações do tipo Van der Waals, presente entre as cadeias poliméricas.
Como respostas medianas, alguns estudantes : • mencionaram o tipo de ligação química como fator principal da dureza diferenciada das diferentes classes de materiais, porém falharam na hierarquização da dureza. • afirmaram que a razão para os polímeros exibirem baixa dureza se deve à presença de ligações secundárias sem identificá -las.
Como exemplo de respostas adequadas, os estudantes : • souberam apontar o tipo de ligação química existente nas diferentes classes de materiais e, a partir daí, definiram a correta hierarquização da dureza, em ordem decrescente, cerâmicos, metais e polímeros, e então associaram corretamente a força de cada ligação, em cada material, justificando tais diferenças de dureza.
De modo geral, nesta questão ocorreu um maior índice de acertos e menor quantidade de questões em branco, quando comparada com as duas específicas anteriores.
Ressalta -se este fato como ponto positivo, já que este tema acessa a ciência dos materiais na sua essência.
A média obtida foi melhor e foram encontradas 217 questões em branco num total de 920 provas, número inferior ao encontrado nas questões anteriores. 3.3.3 Considerações Finais O desempenho demonstrado pelos alunos das Engenharia – Grupo V, no ENADE/2011, relativo às questões discursivas, mostrou uma média global muito baixa.
Este resultado revela a necessidade de melhoria na formação dos alunos.

% matched lemmas: encerramento, reparo, sensato
\end{textsample}
